% 池式主动选择流程；虚线路径保留无法判断状态，不把查询数等同于有效标签数。
\begin{tikzpicture}[x=1mm,y=1mm,font=\figurefont\footnotesize,text=BookInk,
  box/.style={draw=BookRule,fill=BookPaper,line width=.8pt,rounded corners=1.2mm,minimum width=42mm,minimum height=11mm,align=center,inner sep=2.5pt},
  decision/.style={box,draw=FigureModel,fill=FigureModel!10},
  flow/.style={-{Latex[length=2.2mm,width=1.5mm]},draw=BookTeal,line width=1.1pt,rounded corners=1.2mm},
  secondary/.style={-{Latex[length=2.2mm,width=1.5mm]},draw=BookAmber,line width=.85pt,dashed,rounded corners=1.2mm},
  edge label/.style={font=\figurefont\scriptsize,fill=white,inner sep=1.5pt}]
  \node[box,draw=FigureInput,fill=FigureInput!10] (state) at (42,0) {已有知识$L_t$与候选池$U_t$};
  \node[box] (criterion) at (42,-18) {按当前依据计算查询分数};
  \node[box] (select) at (42,-36) {选择样本$x_t$并发出任务};
  \node[decision] (judge) at (42,-54) {人工能否依据现有信息作答？};
  \node[box,draw=FigureOutput,fill=FigureOutput!7] (usable) at (42,-72) {保存答案并加入$L_{t+1}$};
  \node[box,draw=FigureLoss,fill=FigureLoss!10,minimum width=34mm] (pending) at (88,-72) {记录无法判断原因\\不加入已标集合};
  \node[box] (update) at (42,-90) {更新候选池、模型与累计成本};
  \node[decision] (stop) at (42,-108) {预算耗尽或达到完成条件？};
  \node[box,draw=FigureOutput,fill=FigureOutput!7,minimum width=31mm] (end) at (91,-108) {停止并保留未决状态};
  \draw[flow] (state)--(criterion);
  \draw[flow] (criterion)--(select);
  \draw[flow] (select)--(judge);
  \draw[flow] (judge)--node[edge label,right] {可用答案}(usable);
  \draw[secondary] (judge.east)-|node[edge label,pos=.25,above] {无法判断}(pending.north);
  \draw[flow] (usable)--(update);
  \draw[secondary] (pending.south)|-(update.east);
  \draw[flow] (update)--(stop);
  \draw[flow] (stop)--node[edge label,above] {是}(end);
  \draw[flow] (stop.west)--node[edge label,above,pos=.08] {否}(7,-108)|-(criterion.west);
  \node[anchor=west,text=BookMuted] at (66,-93) {$L$只因可用答案增长};
\end{tikzpicture}
